```latex
% Replace the current Introduction with the following.

\section{Introduction: From ``It's the Vibe'' to Vibe Research}

In the Australian film \emph{The Castle}, the inexperienced lawyer Dennis Denuto is convinced that his clients have a valid constitutional case, but is unable to express the underlying principle in formal legal terms. His argument eventually collapses into the memorable phrase, ``It's the vibe.'' At first, the audience is invited to laugh: whatever else legal reasoning may consist of, this surely cannot be enough. But the joke is ultimately on us. Denuto's intuition is substantially correct. When an experienced barrister reformulates it in language recognised by the court, the case succeeds.

The expression \emph{vibe coding} began with something of the same comic quality. Introduced by Andrej Karpathy in February 2025, the term deliberately described a style of programming which sounded like the opposite of serious software engineering: delegate the implementation to an AI, concentrate on describing the desired result, and increasingly stop worrying about the code itself \citep{fowler2026vibecoding}.

Yet the joke changed remarkably quickly. Vibe coding is now described in the research literature as a widespread or emerging software-development paradigm \citep{meskeEtAl2025vibecoding,sarkarDrosos2025vibecoding}, while some practitioner accounts already describe it as the dominant paradigm for interaction with agentic development environments \citep{nimblabs2026vibecoding}. The terminology remains informal and contested, but treating the underlying practice as merely a joke is no longer a plausible description of contemporary software engineering.

In a modern vibe-coding workflow, a human developer may formulate requirements, inspect visible behaviour, request modifications, perform selected independent checks, and ultimately decide whether a system is ready to ship. An advanced coding assistant may meanwhile perform much of the detailed implementation, debugging, refactoring, test construction, repository inspection, and technical explanation.

The point is not that such development is risk-free. It plainly is not. A developer who delegates substantial work to an AI system must decide how much independent checking is appropriate. With experience, this becomes partly a problem of calibrated trust: learning which outputs require close inspection, what kinds of tests provide useful evidence, which failures are characteristic of the system being used, and when additional verification is needed.

One of the principal attractions of contemporary vibe coding is precisely that sufficiently capable coding assistants make exhaustive line-by-line human inspection unnecessary in many cases. If every generated component had to be independently reconstructed or checked in detail, much of the productivity advantage would disappear.

Nor does the supervising human necessarily acquire detailed knowledge of everything the coding assistant has done. Expertise may be distributed between the human and the AI system. If another team member asks why an obscure component was implemented in a particular way, the most useful response may be to ask the coding assistant to inspect the repository and explain it.

These features suggest a natural question. What happens when essentially the same methodology is applied to academic research?

Advanced AI systems can now participate in formulating research questions, reviewing literature, proposing hypotheses, designing experiments, implementing analyses, interpreting results, constructing arguments, drafting papers, and responding to criticism. We use the term \emph{vibe research} for sustained iterative research collaboration of this kind.

Gunkel \citep{gunkel2025prompted} approaches human--AI collaboration from a very different direction. Drawing on Barthes, Derrida, and the history of distributed creative production, he challenges both the instrumental conception of generative AI as a passive tool and the romantic conception of the author as a singular origin of meaning; the essay is itself presented as the product of an iterative collaboration with ChatGPT-4o. Our starting point is instead software-engineering practice, in particular the delegation of substantive work to agents operating within versioned repositories. It is striking that these philosophical and technical routes lead to a related conclusion: the model of a sovereign human author employing a passive instrument describes the observed process poorly, and intellectual production is better understood as distributed, iterative, and dialogical.

The central puzzle of this paper follows directly. If vibe coding and vibe research involve closely similar forms of human--AI collaboration, why has vibe coding become an increasingly accepted development practice while vibe research remains difficult to accommodate within mainstream academic norms?

The history of vibe coding gives some reason to be cautious about dismissing this question on intuitive grounds. To someone unfamiliar with current agentic workflows, the suggestion that an AI could perform most of the implementation work on a serious software project may itself still sound faintly absurd. Very recently, that reaction would have been entirely conventional. Practice changed faster than intuition.

The same possibility deserves consideration in research. The idea that an AI system could be a substantive research collaborator---still more, that it might appropriately be represented as an author---is often treated as inherently implausible. But the recent history of vibe coding suggests that the fact that a new practice initially sounds unserious is not evidence that it will remain so.

Arguably the most important difference is sociological. Research does not consist only in producing intellectual artefacts. Being a researcher normally means participating in a research community: publishing work under one's name, responding to questions and criticism, explaining how results were obtained, and accepting various forms of responsibility for what one has helped produce.

This is often phrased as some version of the question \emph{``Can the AI be a responsible author?''}, with the imputation that it cannot. In current mainstream academic practice, AI authors are generally not permitted. The International Committee of Medical Journal Editors, for example, states that chatbots should not be listed as authors because they cannot be responsible for the accuracy, integrity, and originality of the work \citep{icmje2026ai}. Nature Portfolio similarly states that authorship carries accountability which cannot effectively be applied to large language models \citep{nature2026ai}. The ACL Policy on Publication Ethics explicitly states that generative AI tools may not be listed as authors and that ACL does not consider a generative model capable of fulfilling co-authorship requirements \citep{acl2026policy}.

These policies identify a real issue, but the apparently straightforward concept of ``responsibility'' turns out to cover several different functions. A person may be institutionally responsible for approving a piece of work without possessing all of the detailed expertise embodied in it. Conversely, the person or system best able to explain a particular intellectual decision may have no formal authority over publication.

For the remainder of the paper we therefore avoid treating responsibility as an unanalyzed property. We shift instead to the closely related but more tractable concept of accountability, distinguishing institutional accountability from intellectual accountability.

This distinction leads to the more operational notion of \emph{reliable intellectual provenance}: the ability to establish where important ideas and decisions came from, inspect the evidence on which they were based, and obtain an explanation of them. The resulting practical question is whether a research project involving substantial AI participation can actually be organised so that this kind of provenance is available. Our central technical observation is that the requirements are remarkably similar to those already encountered in serious AI-assisted software development.

But if the technical analogy is as close as we suggest, the original question remains. Why have the two communities developed such different norms? We will argue that part of the explanation may lie not in intrinsic differences between software and research, but in the mechanisms through which their working practices change. Software methodology can evolve directly through individual, team, and organisational practice. Academic research is additionally constrained by formal systems of publication governance whose categories may change on a substantially slower timescale.

Our aim, then, is not to ask the reader to accept ``the vibe'' as an argument. It is to unpack the intuition into a more explicit account of distributed expertise, accountability, provenance, agentic infrastructure, and governance. In that limited sense, we are attempting for \emph{vibe research} something analogous to what the experienced barrister does for Denuto's intuition: to identify more precisely what the practice consists of, why it can work, and how it relates to established institutional concepts.

After developing these arguments, we examine documented examples from our own work and recent developments elsewhere. The present paper is itself being produced through the same methodology, and its development history forms part of the accompanying provenance record.


% Insert the following new subsection after "What Do the Examples Tell Us?"
% and before "Two Timescales of Change: Practice and Governance".

\subsection{Beyond Our Examples: AI as a Source of New Mathematical Knowledge}

The examples above concern our own work and therefore cannot by themselves establish how far substantive AI contribution extends elsewhere. Recent developments in mathematics provide useful independent evidence.

The relevant point is narrower than the claim that AI systems are autonomous mathematicians, or that human mathematicians have become unnecessary. It is that substantive mathematical ideas and results can now originate in AI-mediated processes in which the decisive intellectual content was not supplied in advance by a human researcher.

A comparatively clean early example is Erd\H{o}s Problem~728. Sothanaphan \citep{sothanaphan2026erdos728} describes its resolution as the first Erd\H{o}s problem regarded as fully resolved autonomously by an AI system. The system combined GPT-5.2 Pro with the Aristotle theorem prover and produced a formal Lean proof, subsequently translated into conventional mathematical exposition.

More strikingly, in May 2026 an internal OpenAI model found a counterexample to a long-standing conjecture of Erd\H{o}s concerning the planar unit-distance problem. A group of leading mathematicians subsequently produced a human-verified exposition, explicitly describing the construction as an ``OpenAI-generated counterexample'' \citep{alonEtAl2026unitdistance}. The result did not merely automate a proof whose strategy had been supplied by a human; it produced a construction which overturned a conjecture that had guided work on the problem for decades.

A further case appeared in July 2026, when Levent Alp\"oge announced a counterexample, found using Claude Fable~5, to the Jacobian Conjecture in dimension three. The mathematical counterexample was rapidly checked and has since been independently analysed and generalised \citep{gao2026jacobian,lee2026jacobian}. Here some caution about provenance is appropriate: the publicly available mathematical literature establishes the counterexample much more firmly than it establishes a detailed reconstruction of the respective contributions made by the human and AI participants. For the present argument, however, the episode is further evidence that frontier AI systems are participating in research at a level difficult to characterise as mere text generation.

The most recent example is in analytic number theory. In August 2026, a research version of Claude produced a proof that more than two thirds of the non-trivial zeros of the Riemann zeta function lie on the critical line, substantially improving the previous unconditional lower bound. Anthropic released a complete Lean formalisation of the result, and an independent group subsequently rebuilt and checked the formalisation under separate proof kernels \citep{claude2026zeta,riemannlabs2026zeta}. This is emphatically not a proof of the Riemann Hypothesis, but it is a non-trivial new result on one of the most intensively studied objects in mathematics.

These examples differ greatly in their degree of human involvement, methods of verification, and mathematical significance. They should not be collapsed into a single claim about ``AI doing mathematics''. What they collectively establish is more limited and more relevant here: it is no longer safe to assume, as part of the definition of AI-assisted intellectual work, that substantive ideas necessarily originate with the human participant.

This matters for the present discussion because many intuitive accounts of AI authorship tacitly rely on exactly that assumption. If the AI is understood as a sophisticated instrument which merely executes human intellectual intentions, then excluding it from the provenance record is straightforward. Once AI systems can themselves supply conjectures, constructions, proof strategies, counterexamples, analyses, and explanations, the descriptive problem becomes harder.

The issue is still not settled by saying that the AI made a substantive contribution. Questions of verification, attribution, accountability, and authorship remain. But those questions must now be asked in response to an empirical fact: in at least some current research workflows, intellectually significant content is being generated by the AI participant rather than merely transmitted through it.


% In Discussion and Conclusion, after the paragraph summarising our three examples,
% insert the following:

Recent developments in mathematics provide independent evidence for the same broader point. AI systems have now produced autonomously generated solutions to Erd\H{o}s problems, a counterexample overturning a long-standing conjecture in discrete geometry, substantial contributions to the recent Jacobian Conjecture counterexample, and a formally verified new bound concerning zeros of the Riemann zeta function \citep{sothanaphan2026erdos728,alonEtAl2026unitdistance,gao2026jacobian,claude2026zeta}. These cases differ substantially and should not be treated as equivalent. They nevertheless make one assumption increasingly difficult to sustain: that substantive intellectual content in AI-assisted research must, by definition, originate with the human participant.

% Near the end of Discussion and Conclusion, replace the final two paragraphs with:

We suspect that the practices described here will already be familiar to many researchers. That is an empirical conjecture, not a conclusion.

If it is correct, however, the main novelty of vibe research may turn out not to be the practice itself. It may simply be giving an increasingly common practice a name, documenting it explicitly, and asking whether the concepts and governance structures of the research community still describe accurately how some research is actually being done.

The phrase itself is intentionally informal. So was ``vibe coding'', and so, in a different way, was Dennis Denuto's argument. In \emph{The Castle}, the initial temptation to laugh turns out to be premature: the intuition is not sufficient, but neither is it wrong. It needs to be translated into concepts which the relevant institution can evaluate.

That is the modest ambition of this paper. The question is not whether ``the vibe'' should substitute for an account. It is whether, once the account is supplied, the underlying practice looks quite as absurd as it first appeared.
```
